\section{Related Work}\label{sec:related}
% src: S-PAPER-1
We build on two references from our proposal.
Snowberg and Wolfers study horse racing \citep{snowberg2010favorite}.
They show that longshots win less often than their odds say.
We test the same gap on Polymarket.

The dataset paper reports Polymarket results on cheap tokens \citep[Table~3]{qin2026polymarketv1}.
Tokens priced under \PaperCheapPrice{} pay out \PaperCheapGap{} less than their price.
This is counted by trade.
We count markets, not trades.
We read this result as the same direction as the horse-racing bias.

The dataset paper also scores \PaperBrierMarkets{} markets with a squared-error score \citep{qin2026polymarketv1}.
We score all resolved Yes/No markets that pass our filters.
Finally, the paper reports that the top \TopOneLabel{} of maker addresses hold \PaperTopShare{}
of maker volume \citep[Table~5]{qin2026polymarketv1}.
We ask whether a few wallets dominate volume in our kept markets.
% src: S-PAPER-11
Five further references place our tests in context.
Wolfers and Zitzewitz review prediction markets \citep{wolfers2004prediction}.
They report that market forecasts are typically fairly accurate and usually beat moderately sophisticated benchmarks.
We test this accuracy as calibration of Polymarket prices.
Page and Clemen find that prediction markets are reasonably well calibrated when time to expiration is short \citep{page2013prediction}.
They find prices significantly biased for events farther in the future, with a favourite/longshot bias.
Our longshot check looks at the same bias.
Le uses 353 million trades on 429{,}000 binary contracts from Kalshi and Polymarket \citep{le2026decomposing}.
The preprint finds that calibration varies by event domain, time to resolution, and trade size.
We test all kept markets together and do not split by domain or trade size.

Gneiting and Raftery name sharpness subject to calibration as the goal of probabilistic forecasting \citep{gneiting2007strictly}.
They list the Brier score among the proper scoring rules.
Our mean squared residual of the price is this score.
Clauset, Shalizi, and Newman fit power laws by maximum likelihood \citep{clauset2009power}.
They test the fit with the Kolmogorov-Smirnov statistic and with likelihood-ratio tests against alternatives.
On 24 real datasets, the power law fits some and is ruled out for others.
Like them, we measure the fit with the KS statistic.
We fix the lower bound at the top \TopTenLabel{} of values and run no likelihood-ratio test.
\todo{owner: confirm the added references and add any others}
